\section{Proofs for Section~\ref{sec:readers}: The Reduction to Holes}
\label{app:reduction}

\subsection{Leaf Cones}\label{app:reduction-cones}

We use the following property of trees: for a tree $t$ and a path $w$, at
most one prefix of $w$ is a leaf of $t$, and if none is, then $w$ is a
position of $t$ labeled $f$. Indeed, by induction on its prefixes, a prefix
of $w$ is a position as long as no shorter prefix is a leaf, since a
position that is not a leaf has all its children, and a leaf has none.

\begin{proof}[Proof of \cref{thm:cones}]
\eqref{it:cones-pair} The cones are extension closed by definition, and
disjoint because a path has at most one prefix that is a leaf. Conversely,
let $L^{\sgp},L^{\sgm}$ be disjoint extension-closed sets of paths. Let $D$
be the set of paths none of whose proper prefixes lies in $L^{\sgp}\cup
L^{\sgm}$, and define $t$ on $D$ as follows: $t(w)=f$ if $w\notin
L^{\sgp}\cup L^{\sgm}$; if $w\in L^s$, then $t(w)$ is the label of a leaf of
sign $s$ at polarity $\pol(w)$, that is, $\Top$ if $s=\sgp$ and $\pol(w)=0$
or $s=\sgm$ and $\pol(w)=1$, and $\Bot$ otherwise. The set $D$ contains
$\eps$ and is closed under prefixes; a position labeled $f$ has all its
children in $D$, and a position in $L^{\sgp}\cup L^{\sgm}$ has none. So $t$
is a tree, and its leaves of sign $s$ are the elements of $L^s$ without a
proper prefix in $L^{\sgp}\cup L^{\sgm}$. A path $w$ lies in $L^s$ exactly
when its shortest prefix in $L^{\sgp}\cup L^{\sgm}$ exists and lies in
$L^s$, because both sets are extension closed and they are disjoint; hence
$\Cone st=L^s$. A tree is determined by its cones: its leaves are the
minimal elements of the union of the cones, their labels are determined by
their signs and polarities, and its other positions are the paths without a
prefix in the union.

\eqref{it:cones-order} Let $t\le u$ and $w\in\Cone{\sgp}t$, with a prefix
$v$ of $w$ that is a leaf of $t$ of sign $\sgp$. Let $v'$ be the longest
prefix of $v$ that is a position of $u$; it is a position of both trees. If
$v'=v$ and $\pol(v)=0$, then $t(v)=\Top\le u(v)$, so $u(v)=\Top$; if
$\pol(v)=1$, then $u(v)\le t(v)=\Bot$. In both cases $v$ is a leaf of $u$ of
sign $\sgp$. If $v'$ is a proper prefix of $v$, then $u(v')$ is a leaf and
$t(v')=f$; for $\pol(v')=0$ we get $f\le u(v')$, so $u(v')=\Top$, and for
$\pol(v')=1$ we get $u(v')\le f$, so $u(v')=\Bot$. Again $v'$ is a leaf of
$u$ of sign $\sgp$, and a prefix of $w$. Hence
$\Cone{\sgp}t\subseteq\Cone{\sgp}u$, and the inclusion
$\Cone{\sgm}u\subseteq\Cone{\sgm}t$ follows in the same way with the roles
of the trees and of the labels exchanged.

Conversely, assume both inclusions and let $\pi$ be a position of both
trees with $\pol(\pi)=0$; we show $t(\pi)\le u(\pi)$. If $t(\pi)=\Top$, then
$\pi\in\Cone{\sgp}t\subseteq\Cone{\sgp}u$, so a prefix of $\pi$ is a leaf of
$u$ of sign $\sgp$; since $\pi$ is a position of $u$ this prefix is $\pi$,
and $u(\pi)=\Top$. If $t(\pi)=f$ and $u(\pi)=\Bot$, then
$\pi\in\Cone{\sgm}u\subseteq\Cone{\sgm}t$, so a prefix of $\pi$ is a leaf of
$t$, contradicting $t(\pi)=f$. If $t(\pi)=\Bot$ there is nothing to show.
The case $\pol(\pi)=1$ is symmetric.

\eqref{it:cones-node} The only prefix of $\eps$ is $\eps$, which has
polarity $0$. For $t=f(t_1,\dots,t_n)$, the leaves of $t$ are the paths
$iv$ with $v$ a leaf of $t_i$, with the same label and with
$\pol(iv)=c(i)+\pol(v)$; so the sign of $iv$ in $t$ is the sign of $v$ in
$t_i$ if $i$ is covariant and the opposite sign if $i$ is contravariant.
\end{proof}

\begin{proof}[Proof of \cref{lem:clause-form}]
By \cref{thm:cones}\eqref{it:cones-pair}, $\rho\mapsto S_\rho$ is a
bijection from the assignments onto the consistent sets of atoms that are
closed under extensions. It remains to show that $\rho$ satisfies a literal
if and only if $S_\rho$ is closed under its moves, contains its facts and
contains none of its forbidden atoms. For $x\le y$ this is
\cref{thm:cones}\eqref{it:cones-order}, and for $x=\Bot$ and $x=\Top$ it is
the first statement of \eqref{it:cones-node}. Consider $f(x_1,\dots,x_n)\le
y$ and let $t=f(\rho(x_1),\dots,\rho(x_n))$. By \eqref{it:cones-node},
$\eps$ lies in no cone of $t$, and $iw\in\Cone{\sgp}t$ if and only if
$(x_i^{[i]},w)\in S_\rho$, and $iw\in\Cone{\sgm}t$ if and only if
$(\flip{x_i^{[i]}},w)\in S_\rho$. So $\Cone{\sgp}t\subseteq
\Cone{\sgp}{\rho(y)}$ says that $S_\rho$ is closed under the moves
$(x_i^{[i]},w)\to(y^{\sgp},iw)$, and $\Cone{\sgm}{\rho(y)}\subseteq
\Cone{\sgm}t$ says that $(y^{\sgm},\eps)\notin S_\rho$ and that $S_\rho$ is
closed under the moves $(y^{\sgm},iw)\to(\flip{x_i^{[i]}},w)$. The literal
$y\le f(x_1,\dots,x_n)$ is symmetric.
\end{proof}

\subsection{Derivations}\label{app:reduction-derivations}

A move changes the signed variable of an atom and replaces a prefix of
length at most one of its path by a prefix of length at most one; the set of
moves does not depend on the rest of the path. We record this as two
\emph{stack properties}, for all atoms and all paths $\sigma$:
\begin{enumerate}[label=\textup{(M\arabic*)}]
\item\label{it:M1} if $(p,w)\mv(q,w')$, then $(p,w\sigma)\mv(q,w'\sigma)$;
\item\label{it:M2} if $(p,w\sigma)\mv(q,w'\sigma)$ by a derivation in which
  every atom has a path of the form $z\sigma$, where the displayed
  occurrence of $\sigma$ is never popped, then $(p,w)\mv(q,w')$.
\end{enumerate}
In a derivation by moves, let $\ell$ be the path of an atom of minimal
length. Then every path of the derivation has the form $z\ell$ with this
occurrence of $\ell$ never popped, because a letter of the path can change
only when all letters above it have been popped. We call $\ell$ a
\emph{lowest point} of the derivation.

\begin{proof}[Proof of \cref{lem:duality}]
\eqref{it:dual} is read off \cref{tab:clauses}: the two moves of a literal
$x\le y$ are dual to each other, and so are the two moves of each
constructor literal, for every $i$ and $w$. For \eqref{it:ext-first}, the
implication from right to left holds because $(p,u)\der(p,u\sigma)$ by
extensions. For the other implication we show, by induction on the length of
a derivation of $(q,v)$ from $(p,u)$, that $(p,u\sigma)\mv(q,v)$ where
$\sigma$ consists of the letters of its extension steps in order. If the
last step is a move, this follows from the induction hypothesis. If it is an
extension $(q,v')\to(q,v'i)$, the induction hypothesis gives
$(p,u\sigma')\mv(q,v')$, and \ref{it:M1} gives $(p,u\sigma'i)\mv(q,v'i)$.
\end{proof}

\begin{proof}[Proof of \cref{lem:witness}]
By \cref{thm:cones}\eqref{it:cones-order}, $\rho(x)\not\le\rho(y)$ if and
only if some path lies in $\Cone{\sgp}{\rho(x)}\setminus\Cone{\sgp}{\rho(y)}$
or in $\Cone{\sgm}{\rho(y)}\setminus\Cone{\sgm}{\rho(x)}$.
\end{proof}

\subparagraph*{Runs.}
For a set $P$ of signed variables let $\operatorname{cl}^{\uparrow}(P)$ be
the set of $r$ with $(p,\eps)\mv(r,\eps)$ for some $p\in P$, and
$\operatorname{cl}^{\downarrow}(P)$ the set of $r$ with $(r,\eps)\mv(p,\eps)$
for some $p\in P$. Define
\begin{align*}
  \stepU(P,i)&=\operatorname{cl}^{\uparrow}\bigl(\{\,r'\mid(r,i)\to(r',\eps)
     \text{ is a move for some }r\in P\,\}\bigr),\\
  \stepD(P,i)&=\operatorname{cl}^{\downarrow}\bigl(\{\,r'\mid(r',\eps)\to(r,i)
     \text{ is a move for some }r\in P\,\}\bigr).
\end{align*}
The relation $(p,\eps)\mv(r,\eps)$ between signed variables is computable:
it is the least reflexive and transitive relation $R$ that contains the
moves between atoms with empty paths and satisfies: if $(p,\eps)\to(p',i)$
and $(r',i)\to(r,\eps)$ are moves and $p'\mathrel Rr'$, then $p\mathrel Rr$.
This follows by induction on the length of a derivation from $(p,\eps)$ to
$(r,\eps)$: either some intermediate atom has an empty path, and the
derivation splits, or it has one step, or its first step pushes a letter $i$
that its last step pops, and \ref{it:M2} applies to the part in between.

\begin{proof}[Proof of \cref{lem:runs-det}]
Let $r\in\Up(p,w)$ and let $(r,i)\to(r',\eps)$ be a move and
$(r',\eps)\mv(r'',\eps)$. By \ref{it:M1}, $(p,wi)\mv(r,i)\to(r',\eps)\mv
(r'',\eps)$, so $r''\in\Up(p,wi)$. Conversely, consider a derivation of
$(r'',\eps)$ from $(p,wi)$ and the first atom with an empty path in it. The
atom before it is $(r,j)$ for a letter $j$ and a move $(r,j)\to(r',\eps)$.
All earlier atoms have nonempty paths, so their last letter is never popped
and is $i$; hence $j=i$, and $(p,w)\mv(r,\eps)$ by \ref{it:M2}. The rest of
the derivation shows $r''\in\operatorname{cl}^{\uparrow}(\{r'\})$. This
proves $\Up(p,wi)=\stepU(\Up(p,w),i)$, and $\Up(p,\eps)=
\operatorname{cl}^{\uparrow}(\{p\})$. The statement for $\Dn$ follows in the
same way from the last atom with an empty path in a derivation of $(p,wi)$
from $(r'',\eps)$, or from the first statement by duality. Finally, by
\cref{tab:clauses} a move $(r,i)\to(r',\eps)$ exists only for $r=y^{\sgp}$
with a literal $y\le f(\dots)$ or $r=y^{\sgm}$ with a literal $f(\dots)\le
y$, and then $r\in\NegV$.
\end{proof}

Note that $\stepU(\varnothing,i)=\varnothing$, so $\Up(p,u)=\varnothing$
implies $\Up(p,uz)=\varnothing$ for every $z$. By
\cref{lem:duality}\eqref{it:dual},
\begin{equation}\label{eq:run-dual}
  \Dn(p,w)=\{\,\flip r\mid r\in\Up(\flip p,w)\,\}.
\end{equation}

\begin{proof}[Proof of \cref{lem:cuts}]
Let $a,b$ be cuts as in the statement, $r\in\Up(p,\wpre ua)\cap
\Dn(q,\wpre vb)$ and $\wsuf vb=\wsuf ua\,\sigma$. Then, by \ref{it:M1},
\[
  (p,u)\der(p,u\sigma)=(p,\wpre ua\,\wsuf vb)\mv(r,\wsuf vb)\mv
  (q,\wpre vb\,\wsuf vb)=(q,v).
\]
Conversely, let $(p,u)\der(q,v)$. By \cref{lem:duality}\eqref{it:ext-first},
$(p,u\sigma)\mv(q,v)$ for some $\sigma$. Let $\ell$ be a lowest point of
this derivation, reached at an atom $(r,\ell)$, and write $u\sigma=u'\ell$
and $v=v'\ell$. By \ref{it:M2}, $r\in\Up(p,u')\cap\Dn(q,v')$. If
$\len{u'}>\len u$, then $u'=ujz$ for a letter $j$ and a path $z$, and
$\Up(p,u')\ne\varnothing$ implies $\Up(p,uj)\ne\varnothing$, hence
$\Up(p,u)\cap\NegV\ne\varnothing$ by \cref{lem:runs-det}, contrary to the
hypothesis. So $\len{u'}\le\len u$, and $a=\len{u'}$, $b=\len{v'}$ are cuts
with $\wsuf ua\,\sigma=\ell=\wsuf vb$.

For the last statement, if $r\in\Up(p,u)\cap\NegV$ then $(p,u)\mv(r,\eps)$,
a forbidden atom. Conversely, if $(p,u\sigma)\mv(r,\eps)$ with $r\in\NegV$,
then either $\sigma=\eps$ and $r\in\Up(p,u)$, or $\Up(p,u\sigma)\ne
\varnothing$ with $\sigma\ne\eps$, which implies
$\Up(p,u)\cap\NegV\ne\varnothing$ as above.
\end{proof}

\subsection{The Least Closure and Satisfiability}\label{app:reduction-lc}

Let $R_F$ be the set of signed variables $r$ with $r\in\Up(p,\sigma)$ for a
fact $(p,\eps)$ and a path $\sigma$. By \cref{lem:runs-det} it is the least
set that contains $\Up(p,\eps)$ for every fact $(p,\eps)$ and is closed
under the maps $\stepU(\cdot,i)$ applied to its singletons; so it is
computable.

\begin{lemma}[Least closure]\label{lem:lc}
An atom $(q,v)$ lies in $\LC$ if and only if $\Dn(q,\wpre vb)\cap R_F\ne
\varnothing$ for some cut $b$ of $v$. If $\Gamma$ is satisfiable, then
$R_F=\FctV$.
\end{lemma}

\begin{proof}
$(q,v)\in\LC$ if and only if $(p,\sigma)\mv(q,v)$ for a fact $(p,\eps)$ and
a path $\sigma$, by \cref{lem:duality}\eqref{it:ext-first}. A lowest point
$\ell$ of such a derivation gives $\sigma=\sigma_1\ell$, $v=v'\ell$ and a
signed variable in $\Up(p,\sigma_1)\cap\Dn(q,v')$, as in the proof of
\cref{lem:cuts}; conversely such a variable gives a derivation with
$\sigma=\sigma_1\wsuf vb$, by \ref{it:M1}. If $\Gamma$ is satisfiable, $\LC$
contains no forbidden atom, so $\Up(p,\eps)\cap\NegV=\varnothing$ for every
fact, and $\Up(p,\sigma_1)=\varnothing$ for $\sigma_1\ne\eps$ by
\cref{lem:runs-det}; hence $R_F$ is the union of the sets $\Up(p,\eps)$ over
the facts, which is $\FctV$.
\end{proof}

\begin{proof}[Proof of \cref{lem:sat-decidable}]
By \cref{sec:readers-atoms}, $\Gamma$ is satisfiable if and only if $\LC$
contains no forbidden atom and no atom together with its opposite. By
\cref{lem:lc}, a forbidden atom $(r,\eps)$ lies in $\LC$ if and only if
$\Dn(r,\eps)$ meets $R_F$. And $\LC$ contains $(x^{\sgp},v)$ and
$(x^{\sgm},v)$ for some $v$ if and only if the product of the two automata
that compute $\Dn(x^{\sgp},\cdot)$ and $\Dn(x^{\sgm},\cdot)$
(\cref{lem:runs-det}), with one bit for each component that records whether
its state has met $R_F$, can reach a state with both bits set. Both
conditions are decidable.
\end{proof}

\subsection{Realization}\label{app:reduction-realization}

\begin{lemma}[Return]\label{lem:return}
Let $\Gamma$ be satisfiable, let $a=(p,u)$ be an atom that derives no
forbidden atom, and let $a\der b$.
\begin{enumerate}
\item\label{it:return-lc} If $\flip b\in\LC$, then $\flip a\in\LC$.
\item\label{it:return-seed} If $a'$ is an atom that derives no forbidden
  atom and $a'\der\flip b$, then $a'\der\flip a$ or $a\der\flip{a'}$.
\end{enumerate}
\end{lemma}

\begin{proof}
Let $b=(q,v)$. By \cref{lem:cuts}, whose hypothesis holds by its last
statement, there are cuts $\alpha$ of $u$ and $\beta$ of $v$, a path
$\sigma$ with $\wsuf v\beta=\wsuf u\alpha\,\sigma$, and $r\in
\Up(p,\wpre u\alpha)\cap\Dn(q,\wpre v\beta)$. By \eqref{eq:run-dual},
$\flip r\in\Up(\flip q,\wpre v\beta)\cap\Dn(\flip p,\wpre u\alpha)$, so by
\ref{it:M1}
\begin{equation}\label{eq:return}
  \flip b=(\flip q,\wpre v\beta\,\wsuf v\beta)\mv(\flip r,\wsuf u\alpha\,
  \sigma)\mv(\flip p,u\sigma).
\end{equation}
We use twice that $\Dn(\flip p,z)=\varnothing$ for every path $z$ of which
$u$ is a proper prefix: otherwise $\Up(p,z)\ne\varnothing$ by
\eqref{eq:run-dual}, and then $\Up(p,u)\cap\NegV\ne\varnothing$ as in the
proof of \cref{lem:cuts}.

\eqref{it:return-lc} Since $\LC$ is closed, \eqref{eq:return} gives
$(\flip p,u\sigma)\in\LC$. By \cref{lem:lc} some prefix $z$ of $u\sigma$ has
$\Dn(\flip p,z)\cap\FctV\ne\varnothing$. By the remark, $z$ is a prefix of
$u$, so $(\flip p,u)\in\LC$ by \cref{lem:lc}.

\eqref{it:return-seed} Let $a'=(p',u')$. By \eqref{eq:return},
$a'\der(\flip p,u\sigma)$, so by \cref{lem:cuts} there are cuts $\alpha'$ of
$u'$ and $\gamma$ of $u\sigma$ with $\wsuf{u'}{\alpha'}\pre
\wsuf{(u\sigma)}\gamma$ and $r'\in\Up(p',\wpre{u'}{\alpha'})\cap
\Dn(\flip p,\wpre{(u\sigma)}\gamma)$. By the remark, $\gamma\le\len u$. The
words $\wsuf{u'}{\alpha'}$ and $\wsuf u\gamma$ are both prefixes of
$\wsuf u\gamma\,\sigma$, so they are comparable. If
$\wsuf{u'}{\alpha'}\pre\wsuf u\gamma$, then $a'\der(\flip p,u)$ by
\cref{lem:cuts}. Otherwise $\wsuf u\gamma\pre\wsuf{u'}{\alpha'}$ and, by
\eqref{eq:run-dual}, $\flip{r'}\in\Up(p,\wpre u\gamma)\cap
\Dn(\flip{p'},\wpre{u'}{\alpha'})$, so $a\der(\flip{p'},u')$ by
\cref{lem:cuts}.
\end{proof}

\begin{proof}[Proof of \cref{thm:realization}]
Let $S$ be a solution set that contains every $A_i$ and no $B_i$. It is
closed and contains the facts, so it contains the hull. Hence no $A_i$
derives a forbidden atom, no $B_i$ lies in $\LC$ or is derived from a seed,
and since $S$ is consistent the same holds for the atoms $\flip{A_i}$. This
proves \ref{it:R1}--\ref{it:R3} and that every such solution set contains
the hull.

Conversely, assume \ref{it:R1}--\ref{it:R3}, and let $H$ be the hull. It is
closed, as a union of closed sets, and contains the facts. It contains no
forbidden atom, by \ref{it:R1} and because $\Gamma$ is satisfiable, and no
$B_i$, by \ref{it:R2} and \ref{it:R3}. Suppose $b,\flip b\in H$. They are
not both in $\LC$, which is consistent. If $A_j\der b$ and $\flip b\in\LC$,
then $\flip{A_j}\in\LC$ by \cref{lem:return}\eqref{it:return-lc},
contradicting \ref{it:R2}. If $A_j\der b$ and $A_i\der\flip b$, then
$A_i\der\flip{A_j}$ or $A_j\der\flip{A_i}$ by
\cref{lem:return}\eqref{it:return-seed}, contradicting \ref{it:R3}. So $H$
is consistent, and it is a solution set.
\end{proof}

\subsection{The Reduction}\label{app:reduction-reduction}

\begin{proof}[Proof of \cref{thm:reduction}]
There are $2^{2\len V}$ sets of signed variables, hence at most
$2^{4\len V}+1$ states. The dead state is absorbing by definition. Let
$(P,Q)$ be rejected, that is, $P\cap\NegV=\varnothing$, and let its
successor on a letter be a state $(P',Q')$. By \cref{lem:runs-det},
$P'=\varnothing$. So $(P',Q')$ is rejected, which is \ref{it:S1}, and it is
the target of no admitted pair, which is \ref{it:S2}. The functions
$\stepU,\stepD$ and the sets $\NegV$, $\FctV$ are computable, so the reader
is.

Let $W$ be a tuple of paths, with seeds $A_i=(p_i,w_i)$ and exclusions
$B_i=(p'_i,w_i)$. By \cref{lem:runs-det} and induction on $a$, the state of
coordinate $i$ at a cut $a$ is
$(\Up(p_i,\wpre{w_i}a),\Dn(p'_i,\wpre{w_i}a))$ if the pairs at all cuts up to
$a$ are states, and $\dead$ otherwise. By \eqref{eq:run-dual} and
\cref{lem:lc}, the run of coordinate $i$ never reaches $\dead$ if and only
if $B_i\notin\LC$ and $\flip{A_i}\notin\LC$. Its endpoint is rejected if and
only if $\Up(p_i,w_i)\cap\NegV=\varnothing$, that is, by \cref{lem:cuts}, if
and only if $A_i$ derives no forbidden atom. So all coordinates are
eligible exactly when \ref{it:R1} and \ref{it:R2} hold. Assume this. A pair
of $\Chan_{ij}(W)$ at cuts $a$ of $w_i$ and $b$ of $w_j$, with
$\wsuf{w_j}b\pre\wsuf{w_i}a$, is admitted if and only if
$\Up(p_j,\wpre{w_j}b)$ meets $\Dn(p'_i,\wpre{w_i}a)$ or
$\Dn(\flip{p_i},\wpre{w_i}a)$. By \cref{lem:cuts}, applied to the seed
$A_j$, such cuts exist if and only if $A_j\der B_i$ or
$A_j\der\flip{A_i}$. So no channel contains an admitted pair exactly when
\ref{it:R3} holds. By \cref{thm:realization}, $W$ is a hole if and only if
some solution set contains the seeds and none of the exclusions. The last
statement of the theorem follows with \cref{lem:clause-form,lem:witness}.
\end{proof}
